\input{./._relpath-to-latexroot.ltx}
\documentclass[\latexroot/\projectname]{subfiles}
\input{\latexroot/@resources/subfile-setup.ltx}

\begin{document}

\section{Introduction}
\whenintegrated{\label{sec:intro}} % integrated doc: SST for labels
\setcounter{page}{0}\pagenumbering{arabic}

In the aftermath of the East Asian financial crisis, a debate emerged about the desirability of floating exchange rate regimes. Contrary to the view that floating exchange rate regimes were beneficial and that the severe economic costs of the crisis were due to a lack of exchange rate adjustment, \cite{CalvoReinhart2002} documented a phenomenon since then known as ``Fear of floating.'' This phenomenon refers to the observation that policymakers in many countries are unwilling to let nominal exchange rates experience large fluctuations, preferring to make use of official reserves to absorb external pressures.

Interest in the topic of optimal exchange rate adjustment in response to a contraction in capital inflows has received renewed attention after the Global Financial Crisis. In Europe, various countries experienced sudden stops and current account reversals, which were associated with severe recessions. Importantly, while many have focused on the experience of countries in the euro area, such as Greece, countries outside the monetary union like Hungary and Latvia also witnessed sharp current account reversals, which were accompanied by large depreciations of the exchange rate.

\cite{VernerGyongyosi2018} document that a crucial channel of transmission of the crisis in Hungary was the one of households' foreign-currency debt, whose real value rose substantially upon the exchange rate depreciation, giving rise to large contractions in consumption for the affected households. Motivated by the evidence of the Hungarian experience, we study a model of a small open economy subject to a capital flow reversal, where the exchange rate is flexible and where households hold debt in foreign currency.

To study the role of household heterogeneity in the transmission of foreign shocks, we develop a quantitative Heterogeneous-Agent New-Keynesian Small Open Model Economy (HANKSOME) that embeds domestically incomplete markets into the leading frameworks for open-economy macroeconomics. The household side follows the standard Bewley-\.{I}mrohoro\u{g}lu-Huggett-Aiyagari model that can capture the rich heterogeneity in household income, consumption and wealth. The production side features New Keynesian elements: nominal rigidities in prices and capital adjustment costs. Finally, the household and domestic production blocks are integrated into the standard open-economy macro framework \citep{ObstfeldRogoff2000,Svensson2000,GaliMonacelli2005}.

Introducing domestically incomplete markets into the standard open-economy macroeconomics framework allows us to capture the redistribution of resources across households in response to shocks. With borrowing constraints and incomplete insurance, households are heterogeneous in their marginal propensities to consume (MPCs). To the extent that shocks redistribute resources toward high or low MPC households, the effects of the shock can be amplified or muted relative to the representative agent model.

We calibrate the model to match salient features of the Hungarian economy in the 2000s. We then use our calibrated HANKSOME framework to study the effect of a sustained expansion in the current account and then an unexpected reversal. We find that in this setting a large exchange rate devaluation that achieves full employment is detrimental from a welfare perspective due to the foreign debt revaluation it entails. Our framework can therefore provide a rationalization for ``fear of floating'' even in the absence of contractionary devaluations.

The paper is organized as follows. Section~\ref{sec:model} presents our small open economy incomplete markets model with nominal rigidities. Section~\ref{sec:heterogeneity} discusses the role of heterogeneity. In Section~\ref{sec:calibration} we discuss how we bring the model to the data. We outline our main experiments in Section~\ref{sec:experiments} and describe our results in Section~\ref{sec:suddenstop}. Section~\ref{sec:mechanism} inspects the mechanism, and Section~\ref{sec:conclusion} concludes.

% Smart bibliography: Only include bibliography if standalone AND has citations
\smartbib

\end{document}
